\documentclass[conference]{IEEEtran}
\IEEEoverridecommandlockouts

\usepackage[T1]{fontenc}
\usepackage[utf8]{inputenc}
\usepackage{cite}
\usepackage{url}
\usepackage{amsmath,amssymb}
\usepackage{graphicx}
\usepackage{booktabs}
\usepackage{array}
\usepackage{multirow}
\usepackage{tabularx}
\usepackage{enumitem}
\usepackage{hyperref}
\usepackage{balance}
\usepackage{xspace}

\hypersetup{hidelinks}

\newcommand{\bertopic}{BERTopic\xspace}
\newcommand{\lda}{LDA\xspace}
\newcommand{\nmf}{NMF\xspace}
\newcommand{\hitl}{HITL\xspace}

\begin{document}

\title{Online Topic Modeling with Human-in-the-Loop Support:\
An Implementation-Centered Report on a BERTopic-Based Monitoring System}

\author{%
\IEEEauthorblockN{Krishna Gajera}
\IEEEauthorblockA{College of Engineering\\
San Diego State University\\
kgajera1820@sdsu.edu}
\and
\IEEEauthorblockN{Professor Name}
\IEEEauthorblockA{College of Engineering\\
San Diego State University\\
xxx@sdsu.edu}
}

\maketitle

\begin{abstract}
Conventional topic modeling treats a corpus as static, yet many practical applications require continuous adaptation to evolving text streams. This paper presents a batch-recurring online topic modeling system that employs \bertopic as the primary operational model, with Latent Dirichlet Allocation (\lda) and Non-negative Matrix Factorization (\nmf) serving as comparative baselines. The central engineering contribution is a \emph{train-and-merge} update strategy: for each incoming time window a fresh \bertopic model is trained on the new batch, then merged with the currently deployed model via a similarity-thresholded integration, thereby preserving accumulated topic structure while admitting emergent themes. The platform further implements multi-signal drift detection---comprising topic prevalence change, embedding-space centroid shift, Jensen--Shannon keyword divergence, and topic birth/death tracking---to quantify model stability across windows. A human-in-the-loop layer supports topic relabeling, merging, version history, and rollback through an audited API. The complete system is orchestrated with Prefect, served through FastAPI, and visualized in a Streamlit dashboard. Evaluation on Twitter and AG~News corpora using coherence ($C_v$), diversity, and silhouette metrics demonstrates that \bertopic consistently achieves higher coherence and diversity than both lexical baselines, while the multi-signal drift framework provides actionable temporal monitoring absent from standard topic modeling workflows.
\end{abstract}

\begin{IEEEkeywords}
BERTopic, topic modeling, online analytics, LDA, NMF, drift detection, human-in-the-loop, dashboard systems
\end{IEEEkeywords}

\section{Introduction}
Topic modeling is often introduced through offline experiments in which a fixed corpus is analyzed once and the resulting topics are treated as static. That setup is useful for learning algorithms, but it does not match how many real analytical systems are expected to operate. In practical settings such as customer-support monitoring, social media analysis, or continuously updated content streams, the main challenge is not simply to fit a topic model once. Instead, the challenge is to repeatedly ingest new text, preserve useful topic structure from prior batches, track distributional changes over time, compare competing modeling approaches, and allow analysts to intervene when automatically generated topics are poorly named or semantically redundant.

This report presents the design and implementation of a batch-recurring online topic monitoring system centered on \bertopic. The project is motivated by the need to combine semantic topic discovery with operational reliability. Specifically, \bertopic is used as the deployed model for topic extraction and inference, while \lda and \nmf are integrated as benchmark baselines to provide comparative evidence under a common evaluation framework.

The core technical idea is a \emph{train-and-merge} temporal update strategy. For each incoming time window, a fresh \bertopic model is trained on the new batch and then merged with the currently deployed model. This strategy allows the system to preserve prior topic structure while incorporating emerging themes, and it aligns naturally with model versioning, rollback, and drift monitoring requirements. Beyond modeling, the implementation includes API-based serving, dashboard analytics, and human-in-the-loop (\hitl) controls for topic relabeling and consolidation.

The main contributions of this project are as follows:
\begin{enumerate}[leftmargin=*]
    \item An end-to-end architecture for online topic intelligence, integrating ingestion, modeling, monitoring, and interaction layers.
    \item A practical \bertopic update mechanism based on batch retraining and model merging for temporal adaptation.
    \item A comparative three-model framework (\bertopic, \lda, \nmf) with shared metrics for defensible model selection.
    \item A multi-signal drift detection module using prevalence change, centroid shift, keyword divergence, and topic birth/death events.
    \item A governed \hitl workflow with audit logs, model version history, and rollback support.
\end{enumerate}

The remainder of this paper is organized as follows. Section II reviews related work. Section III presents system overview and scope. Section IV describes system architecture. Sections V--VIII present processing methodology, model design, and monitoring mechanisms. Section IX discusses evaluation and results. Section X analyzes strengths and limitations, and Section XI concludes the report.

\section{Related Work}
\subsection{Classical Topic Modeling Baselines}
Probabilistic and matrix-factorization approaches remain foundational in topic modeling. \lda models documents as mixtures of latent topics and topics as distributions over words \cite{lda}. \nmf factorizes a document-term matrix into non-negative latent components and often produces interpretable keyword groups \cite{nmf}. These methods are computationally efficient and widely studied, but they are strongly tied to lexical overlap and may degrade on short, noisy, or highly paraphrastic text streams.

\subsection{Semantic Topic Modeling with BERTopic}
Recent methods leverage contextual embeddings to improve semantic grouping quality. \bertopic combines sentence embeddings \cite{sentencetransformers}, dimensionality reduction through UMAP \cite{umap}, density-based clustering with HDBSCAN \cite{hdbscan}, and class-based TF--IDF topic representation \cite{bertopic}. This pipeline is particularly suitable for evolving corpora in which semantically related documents may share limited surface vocabulary.

\subsection{Online Operation, Drift, and Human Oversight}
Most educational topic modeling workflows remain offline and single-pass, whereas deployed systems require temporal updates, monitoring, and analyst control. Practical MLOps patterns emphasize orchestration, reproducibility, and lifecycle observability \cite{prefect,mlflow}. In this context, the present work contributes an implementation-focused integration: batch-recurring model updates via train-and-merge, multi-signal topic drift monitoring, and a human-in-the-loop interface for relabeling, topic consolidation, version tracking, and rollback. This positioning distinguishes the project from standalone modeling experiments by addressing both modeling quality and operational governance.

\section{System Overview and Scope}
This section formalizes the goal of the project and the precise scope of the implemented system. Let a corpus arrive as a sequence of time-indexed batches \(\{\mathcal{B}_t\}_{t=1}^{T}\), where each batch \(\mathcal{B}_t=\{d_i\}\) contains documents observed within a configured time window. The system is designed to (i) extract a human-interpretable topic structure from each window, (ii) maintain an evolving deployed topic model \(\mathcal{M}_t\) across windows, and (iii) provide monitoring signals that quantify when topic structure or usage has shifted sufficiently to require analyst attention. The emphasis is operational: the output is not only a set of topics, but a repeatable lifecycle that supports ingestion, update, evaluation, monitoring, and governance.

\subsection{Operational vs. comparative modeling}
The implementation uses a role-separated modeling strategy. \bertopic is treated as the \emph{operational} model: it is the model that produces the deployed topic catalog, supports inference on unseen text, and drives the dashboard-level analytics. In contrast, \lda and \nmf are implemented as \emph{comparative baselines}. They are executed to contextualize topic quality under a shared evaluation protocol (e.g., coherence, diversity, and clustering separability), but they are not used for serving or human-in-the-loop editing. This distinction is methodologically important because it prevents the evaluation from implicitly assuming that a single modeling family is correct by definition.

\subsection{Temporal update principle: train-and-merge}
The defining design choice for temporal operation is a batch-recurring \emph{train-and-merge} update mechanism. For each new batch \(\mathcal{B}_t\), the system trains a fresh \bertopic model \(\mathcal{F}_t\) using the batch documents, then merges \(\mathcal{F}_t\) with the currently deployed model \(\mathcal{M}_{t-1}\) using a similarity-thresholded integration step. The resulting merged model becomes the new deployed state \(\mathcal{M}_t\). Conceptually, this implements an engineering trade-off: the approach is simpler and more controllable than fully streaming optimization, while still enabling new topics to emerge and old topics to persist across windows. It also aligns with model governance, since each update produces a versionable artifact that can be archived and rolled back.

\subsection{Monitoring and governance objective}
Temporal operation makes monitoring a first-class requirement. The system therefore computes drift signals between consecutive deployed states to detect meaningful change. In the implemented design, drift is not reduced to a single scalar; instead, multiple complementary indicators are tracked, including changes in topic prevalence (distribution shift), movement of topic centroids in embedding space (semantic shift), divergence of topic keyword distributions (descriptor shift), and topic birth/death events (structural shift). These signals support alerting and prioritization in the user interface. In addition, the system exposes human-in-the-loop controls that allow analysts to impose interpretability and consolidation constraints (topic relabeling and merging) while preserving an audit trail and version history.

Figure~\ref{fig:e2e-overview} summarizes the end-to-end system at a conceptual level. It should be read as the operational contract of the platform: each batch window produces updated model state, persisted artifacts, and monitoring outputs that are accessible through API and dashboard layers.

\begin{figure}[t]
\centering
\fbox{\parbox{0.92\columnwidth}{\centering
\vspace{0.75cm}
Placeholder (use your existing end-to-end diagram).\\
Suggested: ingestion $\rightarrow$ preprocessing $\rightarrow$ BERTopic train/merge\\
$\rightarrow$ assignments/metadata $\rightarrow$ LDA/NMF benchmarking $\rightarrow$ drift alerts\\
$\rightarrow$ FastAPI serving $\rightarrow$ Streamlit dashboard $\rightarrow$ HITL relabel/merge/rollback.
\vspace{0.75cm}}}
\caption{End-to-end overview of the implemented online topic monitoring system.}
\label{fig:e2e-overview}
\end{figure}

\subsection{Scope Boundaries}
The system is intentionally scoped to a batch-recurring operational setting. First, updates are performed at window granularity rather than via continuous streaming optimization. Second, \hitl support is strongest for relabeling and merging; topic splitting is not implemented as a fully automated semantic operation. Third, while the software stack is reproducible, complete experiment replay depends on access to the external raw datasets referenced in configuration. These boundaries define the valid interpretation of the experimental evidence presented later.

\section{System Architecture}
The project is implemented as a modular analytics system rather than a monolithic script. At a high level, the architecture has four layers: orchestration, modeling, serving, and interaction.

\subsection{Orchestration layer}
At the orchestration level, the system uses a scheduled or manually triggered batch pipeline. Each run is associated with a time window and a batch identifier. The pipeline determines what data fall inside the current window, processes those documents, updates the primary model, runs benchmarking models if enabled, and stores the outputs required for monitoring and comparison. This design makes the system batch-recurring rather than continuously streaming, which is a practical compromise for a final project: it captures temporal evolution while keeping implementation complexity manageable.

\subsection{Modeling layer}
The modeling layer contains the main \bertopic pipeline and the two comparison baselines. \bertopic functions as the semantic model that powers discovery, metadata generation, and inference. \lda and \nmf are trained on the same batch to provide comparison metrics such as coherence, diversity, and silhouette-style separability. Conceptually, this layer separates \emph{operational use} from \emph{model assessment}. That separation is a strength because it prevents the project from treating a single model as correct by definition.

\subsection{Serving layer}
The serving layer exposes analytical outputs through a web API. This allows the dashboard to request topic summaries, trend histories, benchmark metrics, alerts, and inference results without directly coupling the user interface to the training logic. Such a design is more robust and reviewable than placing all modeling, storage, and visualization in one interactive script.

\subsection{Interaction layer}
The interaction layer is implemented as a dashboard with separate views for overview, current topics, drill-down inspection, drift alerts, inference, benchmarking, and human-in-the-loop operations. From a software engineering perspective, this is important because it turns the project into an analytical tool rather than only a modeling pipeline. A reviewer can understand not just how topics are learned, but also how the system would be used by an analyst.

% TODO: Insert latest architecture diagram here.
\begin{figure}[t]
\centering
\fbox{\parbox{0.92\columnwidth}{\centering
\vspace{0.95cm}
Placeholder for updated architecture diagram.\\
Suggested contents: recurring pipeline, \bertopic as the primary model,\\
\lda and \nmf benchmarking branches, serving API, dashboard, drift monitoring,\\
and human-in-the-loop topic maintenance.
\vspace{0.95cm}}}
\caption{System architecture of the implemented online topic monitoring platform.}
\label{fig:architecture}
\end{figure}

\section{End-to-End Processing Logic}
Although the system is modular, its end-to-end logic follows a clear sequence. Each run begins by identifying the active dataset profile and the time window to be processed. The selected data are then loaded and cleaned according to the active profile. After validation and deduplication, the cleaned texts are passed to the main topic modeling stage.

The main topic modeling stage has two modes. If the system has not yet been initialized for the selected dataset, it performs seed training and produces the first deployable \bertopic model. If a current model already exists, the system trains a fresh \bertopic model on the new batch and merges that model with the existing deployed model. This merged model becomes the new current model. The resulting topics are then used to create metadata, produce document-topic assignments, and update persistent outputs for later dashboard use.

After the primary model update, the system can execute \lda and \nmf on the same batch and record their metrics for comparison. If a prior \bertopic state is available, the system also computes drift signals between the current and previous topic configurations. Finally, run-level metrics and state information are recorded so that subsequent batches can be interpreted in temporal context.

Conceptually, the execution flow can therefore be summarized as:
\begin{enumerate}[leftmargin=*,itemsep=2pt]
    \item identify the current data window;
    \item preprocess and validate the batch;
    \item update the primary \bertopic model;
    \item compute metadata and assignments;
    \item benchmark against \lda and \nmf;
    \item estimate drift against the previous state;
    \item expose results for dashboard and API consumption.
\end{enumerate}

This sequence reflects a monitoring-oriented design. The system is not organized around a single train/test split. Instead, it is organized around repeated operational cycles that allow topic structure to be tracked over time.

\section{Dataset Handling and Preprocessing}
The implementation supports more than one dataset profile, and preprocessing is adapted to the active profile. This is an important design decision because topic modeling quality is highly sensitive to noise patterns in the text. A customer-support corpus has different artifacts from a news corpus, and the preprocessing strategy should reflect that.

In the current implementation, preprocessing includes text normalization, noise removal, empty-text filtering, minimum-length filtering, and duplicate removal. The system distinguishes between at least two cleaning styles: one intended for short, noisy support-style messages and one intended for more general text. The support-oriented cleaner removes artifacts such as URLs, user handles, signature fragments, phone-like strings, repeated punctuation, and emoji-heavy noise. The general cleaner removes more generic formatting noise while avoiding support-specific heuristics that might distort cleaner news-like text.

This dataset-aware preprocessing strategy is a practical strength of the project. It acknowledges that topic modeling quality depends not only on the model but also on preserving the right lexical signal. If preprocessing is too weak, the model learns artifacts; if it is too aggressive, it removes semantic distinctions that are useful for topic formation.

A second practical feature is the generation of document identifiers and batch identifiers. These identifiers are not only administrative. They link raw documents, topic assignments, dashboard summaries, and drift signals across repeated runs. Without such identifiers, it would be difficult to interpret changes over time or to associate user edits with specific model states.

One important limitation, however, is reproducibility from the uploaded archive alone. The code clearly expects external raw data files for the active datasets, but those full datasets are not bundled with the uploaded materials. As a result, the implementation can be analyzed in depth, but some runs cannot be fully replayed end-to-end without the original data sources.

\section{Primary Model Design: BERTopic as the Operational Core}
\subsection{Why \bertopic is the base model}
The project places \bertopic at the center of the system because it is well suited to semantically rich topic discovery on short or moderately noisy texts. Unlike purely count-based topic models, \bertopic begins with sentence embeddings, which allow semantically similar texts to remain close even when they do not share many surface words. Those embeddings are then compressed into a lower-dimensional space, clustered, and converted into interpretable topic descriptors through class-based TF--IDF.

From an implementation perspective, this is a coherent choice for a customer-support or evolving-text setting. Support messages often contain paraphrases, abbreviations, and sparse lexical overlap. A contextual embedding model can help group semantically similar requests that would be fragmented by a purely lexical method.

\subsection{Implemented \bertopic pipeline}
The implemented \bertopic stack follows the standard four-stage pattern: embedding, dimensionality reduction, clustering, and topic representation.
\begin{enumerate}[leftmargin=*,itemsep=2pt]
    \item \textbf{Embedding.} The system uses a sentence-transformer model to produce dense representations of documents.
    \item \textbf{Dimensionality reduction.} A manifold-learning stage reduces the embedding space to a lower dimension that is more tractable for clustering.
    \item \textbf{Clustering.} A density-based clustering method groups nearby reduced points while allowing some items to remain unassigned as noise.
    \item \textbf{Topic representation.} A vectorizer and class-based TF--IDF convert document clusters into interpretable top-word representations.
\end{enumerate}

In the uploaded implementation, the configuration favors relatively small batch settings: neighborhood and clustering thresholds are permissive, and vocabulary filtering is light. This suggests that the project is tuned for moderate batch sizes where preserving topic granularity is valuable, though it also implies sensitivity to noise and parameter choice.

\subsection{Batch-wise online updating}
The most important design choice in the project is the way online updating is implemented. The system does not continuously optimize a single mutable topic model in place. Instead, it uses a batch-wise update scheme. For each new window of data, a fresh \bertopic model is trained on the new batch. That batch model is then merged with the currently deployed model, and the merged artifact becomes the new current model.

This design has several advantages. First, it preserves accumulated topic structure from prior batches while still allowing new themes to enter the system. Second, it fits naturally with versioning, because each update produces a distinct model state that can be archived. Third, it integrates well with human-in-the-loop corrections: if topic labels or merges have already been applied to the current model, those edits can be carried forward into the next merged state.

At the same time, this strategy is not equivalent to fully online optimization in a streaming sense. Topic identity may shift across merges, and numerical topic identifiers do not automatically guarantee semantic continuity. Therefore, while the project is rightly described as ``online'' in the sense of recurring model updates, it is more precisely a \emph{batch-recurring online monitoring system}.

\subsection{Topic metadata and interpretability}
An important implementation detail is the generation of explicit topic metadata after modeling. For each non-outlier topic, the system stores topic identifiers, representative words, estimated size, time-window information, and human-readable labels. It can also augment these topic descriptors with short natural-language summaries. This metadata is essential because the dashboard and user interaction layer do not operate directly on raw cluster assignments. They operate on curated topic descriptions that make the system interpretable.

From a project-review standpoint, this is a meaningful strength. Many student implementations stop at producing topic-word lists. This system goes further by constructing a reusable topic catalog that supports inspection, comparison, and editing.

\section{Comparison Models: \lda and \nmf}
The project includes two benchmark models so that \bertopic is not evaluated in isolation. Both comparison models are implemented in the pipeline and not merely discussed in the documentation.

\subsection{Latent Dirichlet Allocation}
\lda serves as the probabilistic bag-of-words baseline. It models documents as mixtures of latent topics and topics as distributions over words. In the implemented system, \lda is trained on preprocessed lexical features rather than contextual embeddings. This makes it a useful comparison because it represents a classical topic-modeling approach with strong interpretability, but weaker semantic flexibility on short and noisy text.

The presence of \lda in the implementation matters for two reasons. First, it provides a familiar reference model for academic comparison. Second, it helps show whether the observed quality gains from \bertopic are simply due to more aggressive preprocessing or whether they reflect the value of semantic embeddings and clustering.

\subsection{Non-negative Matrix Factorization}
\nmf serves as a second lexical baseline. Unlike \lda, which is probabilistic, \nmf factorizes a non-negative document-term matrix into document-topic and topic-word matrices. In practice, \nmf often yields sharp topic-word lists and can behave differently from \lda even when both are trained on the same corpus.

The implementation contains several stability-oriented decisions in the \nmf path. These indicate that the authors encountered practical difficulties when applying matrix factorization to small and sparse short-text batches. The design therefore prioritizes robust training over strict adherence to user-configured penalties. That is an engineering-driven rather than purely theoretical choice, and it is exactly the kind of detail that should be acknowledged in a final report.

\subsection{Interpretation of the three-model comparison}
Taken together, the three-model setup reflects a sensible hierarchy. \bertopic is used where semantic quality and user-facing interpretability matter most. \lda provides a classical probabilistic comparison. \nmf provides a factorization-based comparison that can reveal different trade-offs in topic sharpness and geometric separation. Because these models operate in different representational spaces, no single metric can declare an absolute winner. Nevertheless, their joint inclusion makes the evaluation more credible than a single-model narrative.

\begin{table}[t]
\caption{Implemented role of each model in the system}
\label{tab:model-role}
\centering
\footnotesize
\begin{tabular}{p{0.18\columnwidth} p{0.23\columnwidth} p{0.47\columnwidth}}
\toprule
\textbf{Model} & \textbf{Role} & \textbf{Purpose in the system} \\
\midrule
\bertopic & Primary model & Topic discovery, online updating, metadata generation, dashboard inspection, inference, and human-guided maintenance \\
\lda & Baseline 1 & Classical probabilistic benchmark for topic quality comparison \\
\nmf & Baseline 2 & Matrix-factorization benchmark for complementary comparison of lexical topic structure \\
\bottomrule
\end{tabular}
\end{table}

\section{Drift Monitoring and Human-in-the-Loop Support}
\subsection{Drift monitoring}
Because the system is designed for repeated batch updates, monitoring topic stability is a necessary part of the architecture. The implementation therefore includes drift assessment between the current and previous topic model states. The drift logic considers multiple signals rather than relying on a single score. These signals include changes in topic prevalence, movement in topic centroids, divergence in topic keywords, and the appearance or disappearance of topics.

This multi-signal design is theoretically appropriate. Topic drift can manifest in different ways: a topic may persist but become much more common, a topic may remain common while its semantics shift, or a new topic may emerge entirely. By monitoring several signals, the system avoids reducing drift to one simplistic indicator.

\subsection{Human-in-the-loop topic maintenance}
The project also supports human-in-the-loop maintenance. In practice, this means the user can inspect topics and intervene when the automatically generated structure is unsatisfactory. The strongest implemented operations are topic relabeling and topic merging. Relabeling improves interpretability by replacing mechanically generated descriptors with analyst-approved names. Merging addresses a common weakness in topic modeling systems: semantically overlapping topics that are distinct numerically but redundant analytically.

The project also includes audit and version history concepts, which are important in a system with manual intervention. Once users begin changing topic labels or merging topics, the system must preserve a record of what changed and when. That versioning capability makes the human-in-the-loop design materially stronger.

A point of caution is that not every desired editing action is equally mature. A topic-splitting interface is present at the system level, but the code indicates that true semantic split functionality is not yet implemented to the same degree as merge and relabel. For an honest final presentation, the system should therefore be described as fully supporting relabeling and merging, with split functionality remaining partial.

% TODO: Insert screenshot of HITL editor or topic management page.
\begin{figure}[t]
\centering
\fbox{\parbox{0.92\columnwidth}{\centering
\vspace{0.8cm}
Placeholder for dashboard screenshot showing topic relabeling, topic merge,\\
or model history / rollback workflow.
\vspace{0.8cm}}}
\caption{Suggested visualization for the human-in-the-loop topic maintenance interface.}
\label{fig:hitl}
\end{figure}

\section{Dashboard and User-Facing Analytics}
A notable strength of the project is that the learned topics are not treated as hidden backend artifacts. The dashboard exposes them through multiple analytical views, allowing a reviewer or analyst to move from a high-level summary to fine-grained topic inspection.

The main dashboard supports at least four important kinds of user-facing analysis. First, it shows the current topic catalog, which makes the active semantic structure of the corpus visible. Second, it provides batch-level or temporal summaries, allowing the user to see how the distribution of topics changes over time. Third, it exposes alerts and drift-related information, connecting model updates to monitoring logic. Fourth, it allows direct inference on unseen text, demonstrating that the deployed \bertopic model is not only a training artifact but also a serving component.

The benchmarking view is especially important in an academic setting because it connects the deployed model to comparative evaluation. Rather than only visualizing the primary model, the system surfaces coherence, diversity, and silhouette-style metrics across \bertopic, \lda, and \nmf. This makes the project easier to defend because the modeling choice is supported by evidence rather than by preference alone.

\section{Experimental Evidence and Benchmark Interpretation}
The uploaded code clearly implements the infrastructure required to compute benchmark histories for \bertopic, \lda, and \nmf. However, the uploaded archive does not include all raw data and generated outputs required to rerun every experiment directly from scratch. For that reason, the report separates \emph{implementation evidence} from \emph{run evidence}. Implementation evidence comes from the source code itself. Run evidence comes from the result image supplied separately by the user.

The supplied benchmark image corresponds to a Twitter dataset experiment and reports batch-wise values for coherence, diversity, silhouette score, and topic count across the three models. The average values visible in that image are summarized in Table~\ref{tab:twitter-results}.

\begin{table}[t]
\caption{Average benchmark values transcribed from the user-supplied Twitter results image}
\label{tab:twitter-results}
\centering
\footnotesize
\begin{tabular}{lccc}
\toprule
\textbf{Model} & \textbf{Coherence} & \textbf{Diversity} & \textbf{Silhouette} \\
\midrule
\bertopic & 0.8126 & 0.9010 & 0.0348 \\
\lda & 0.6503 & 0.2768 & 0.0029 \\
\nmf & 0.4003 & 0.4125 & 0.0924 \\
\bottomrule
\end{tabular}
\end{table}

Several careful observations can be made from these results. First, \bertopic shows the strongest coherence and diversity in the provided experiment. This is consistent with the design of the implementation, which uses contextual embeddings and semantic clustering before deriving interpretable top-word labels. Second, \lda performs substantially worse than \bertopic on both coherence and diversity, which is plausible for short, noisy text where lexical overlap is sparse. Third, \nmf achieves the highest silhouette score in the supplied table, suggesting stronger separation in its feature space even though its topic quality appears weaker than \bertopic according to coherence and diversity.

These results should be interpreted with methodological caution. Silhouette scores are not perfectly symmetric across the three model families because each model assigns topics from a different representational space. Therefore, the strongest evidence in this comparison comes from the joint pattern rather than any one metric in isolation: \bertopic appears to provide the best balance of interpretability-oriented quality measures, while \nmf may produce sharper partitioning in its lexical factorization space.

% TODO: Insert screenshot or recreated plot of the benchmarking page here.
\begin{figure}[t]
\centering
\fbox{\parbox{0.92\columnwidth}{\centering
\vspace{0.8cm}
Placeholder for benchmark chart or screenshot of BERTopic vs. LDA vs. NMF comparison.
\vspace{0.8cm}}}
\caption{Suggested figure for comparative model benchmarking across batches.}
\label{fig:benchmark}
\end{figure}

\section{Implementation Strengths}
The project has several strengths that are especially relevant in a final review.

First, it treats topic modeling as a system rather than as an isolated algorithm. The implementation includes orchestration, monitoring, benchmarking, serving, and user interaction. This makes the project more substantial than a one-off modeling experiment.

Second, the choice of \bertopic as the operational model is well aligned with the type of data the system is intended to analyze. Short support-like texts benefit from semantic representations, and the implementation takes advantage of that while still preserving interpretable topic descriptors.

Third, the inclusion of \lda and \nmf as real implemented baselines strengthens the evaluation. The project does not rely on a single-model argument. Instead, it uses comparison models to justify the modeling choice.

Fourth, the drift-monitoring and human-in-the-loop features show awareness of the practical lifecycle of topic models. In real usage, topics are rarely ``done'' after training. They require supervision, renaming, consolidation, and longitudinal monitoring. The project reflects that reality.

Finally, the dashboard and API design improve accessibility. The system is understandable not only to someone reading code, but also to someone interacting with it as an analytical product.

\section{Limitations and Areas for Refinement}
The project also has important limitations, and these are best presented directly rather than as vague future work.

The first limitation is reproducibility from the uploaded package alone. Although the codebase is substantial, the raw datasets and some generated outputs are not fully included, which means not every experiment can be rerun immediately from the archive.

The second limitation is partial inconsistency in the transition from a support-specific prototype toward a more general multi-dataset system. Some parts of the implementation clearly support multiple dataset profiles, but some surrounding assumptions still reflect an earlier customer-support-centered design. This does not invalidate the system, but it does indicate that the generalization process is not fully complete.

The third limitation concerns human-in-the-loop editing depth. Topic relabeling and merging are meaningfully implemented, but true topic splitting is not yet at the same level of maturity. This matters because split is conceptually harder than merge: it requires discovering coherent substructure within an existing topic rather than only consolidating redundant topics.

The fourth limitation concerns preprocessing consistency. The implementation contains dataset-aware cleaning, which is a strength, but end-to-end consistency must be maintained between training and inference for each dataset profile. Any mismatch there can weaken the reliability of live inference.

The fifth limitation is that model comparison metrics are informative but not perfectly unified across model families. Coherence, diversity, and silhouette each illuminate a different aspect of behavior, but they do not collapse into a single objective notion of quality. This means the evaluation should be presented as comparative evidence rather than as a definitive proof of superiority.

\section{Conclusion}
This project implements a strong final-year style topic modeling system centered on \bertopic and augmented with \lda and \nmf benchmarking, drift monitoring, dashboard-based analytics, and human-in-the-loop maintenance. Its main technical value lies not in proposing a new topic model, but in operationalizing a modern semantic topic model inside a recurring analytical workflow.

The implementation makes several sound design choices. It uses \bertopic as the live model where semantic grouping and interpretability matter most. It benchmarks that model against two classical alternatives. It tracks batch-wise evolution over time rather than treating the corpus as static. It provides mechanisms for monitoring drift and for allowing analysts to correct topic structure. These are meaningful strengths for a project intended to demonstrate both machine learning understanding and systems integration.

At the same time, the system is best described precisely: it is a batch-wise online monitoring system, not a continuously streaming learner; it supports mature relabel and merge operations, but split remains partial; and its results are partly demonstrated through external run artifacts because the uploaded package is not fully self-contained. Presenting the project in this way is not a weakness. On the contrary, it shows technical maturity, because the implementation is explained in terms of what it truly does.

\section*{Acknowledgment}
% TODO: Replace with project/course acknowledgments if desired.
The author may add course, instructor, or team acknowledgments here.

\begin{thebibliography}{99}

\bibitem{bertopic}
M. Grootendorst, ``BERTopic: Neural topic modeling with a class-based TF-IDF procedure,'' \emph{arXiv preprint arXiv:2203.05794}, 2022.

\bibitem{lda}
D. M. Blei, A. Y. Ng, and M. I. Jordan, ``Latent Dirichlet Allocation,'' \emph{Journal of Machine Learning Research}, vol. 3, pp. 993--1022, 2003.

\bibitem{nmf}
D. D. Lee and H. S. Seung, ``Learning the parts of objects by non-negative matrix factorization,'' \emph{Nature}, vol. 401, no. 6755, pp. 788--791, 1999.

\bibitem{sentencetransformers}
N. Reimers and I. Gurevych, ``Sentence-BERT: Sentence embeddings using Siamese BERT-networks,'' in \emph{Proc. EMNLP-IJCNLP}, 2019, pp. 3982--3992.

\bibitem{umap}
L. McInnes, J. Healy, and J. Melville, ``UMAP: Uniform manifold approximation and projection for dimension reduction,'' \emph{arXiv preprint arXiv:1802.03426}, 2018.

\bibitem{hdbscan}
R. J. G. B. Campello, D. Moulavi, and J. Sander, ``Density-based clustering based on hierarchical density estimates,'' in \emph{Advances in Knowledge Discovery and Data Mining}, 2013, pp. 160--172.

\bibitem{prefect}
Prefect Technologies, ``Prefect documentation,'' [Online]. Available: \url{https://www.prefect.io/}. Accessed: Mar. 2026.

\bibitem{mlflow}
M. Zaharia \emph{et al.}, ``Accelerating the machine learning lifecycle with MLflow,'' \emph{IEEE Data Engineering Bulletin}, vol. 41, no. 4, pp. 39--45, 2018.

\end{thebibliography}

\balance
\end{document}
